% Generated draft: requires mathematical and source review.
\documentclass[11pt]{article}
\usepackage[margin=25mm]{geometry}
\usepackage{fontspec}
\setmainfont{Latin Modern Roman}
\usepackage{amsmath,amssymb,mathtools,unicode-math}
\setmathfont{Latin Modern Math}
\usepackage{enumitem,needspace,xurl,hyperref}
\setlength{\parindent}{0pt}
\setlength{\parskip}{5pt}
\setlength{\emergencystretch}{3em}
\newcommand{\N}{\mathbb{N}}
\newcommand{\Z}{\mathbb{Z}}
\newcommand{\Q}{\mathbb{Q}}
\newcommand{\R}{\mathbb{R}}
\newcommand{\C}{\mathbb{C}}
\newcommand{\F}{\mathbb{F}}
\newcommand{\A}{\mathbb{A}}
\newcommand{\PP}{\mathbb{P}}
\newcommand{\E}{\mathbb{E}}
\newcommand{\eps}{\varepsilon}
\newcommand{\dd}{\,\mathrm{d}}
\begin{document}
\section*{UnsolvedMath ID: 30000653}
\textbf{Source problem code:} catalogue entry
\subsection{Definitions and mathematical statement}
\textbf{Commuting derivations conjecture \(\mathrm{CD}(n)\).} Let \(A=\C[x_1,\ldots,x_n]\) and let \(D_1,\ldots,D_{n-1}\) be pairwise commuting locally nilpotent \(\C\)-derivations, linearly independent over \(A\). Must there be a coordinate \(f\in A\) such that\[\bigcap_{i=1}^{n-1}\ker D_i=\C[f]?\] The nontrivial range is \(n\ge2\); for \(n=1\), with the empty-intersection convention, the assertion simply concerns \(A=\C[x_1]\). The source does not explicitly state a dimension range.\par A derivation is \(\C\)-linear and satisfies \(D(ab)=aD(b)+D(a)b\). Locally nilpotent means that for every \(a\in A\) there is an integer \(m\ge1\), allowed to depend on \(a\), with \(D^m(a)=0\). Pairwise commutation means \(D_iD_j=D_jD_i\). Independence over \(A\) means that \(\sum_i a_iD_i=0\) as a derivation, with \(a_i\in A\), forces all \(a_i=0\). It is stronger than independence over \(\C\), but does not assert pointwise independence of the resulting vector fields at every point of affine space.\par A coordinate means that \(f\) extends to a polynomial coordinate system: there are \(f_2,\ldots,f_n\in A\) with \(A=\C[f,f_2,\ldots,f_n]\). The requested conclusion is therefore stronger than saying the common kernel is abstractly a polynomial ring in one variable, or is generated by one arbitrary polynomial. No common slice, everywhere-free action, or stronger independence condition is inserted.\par The supplied triage checked 2026-08-21 labels the conjecture open and reports no located complete resolution, with the report as its sole source. It provides no precise dimension-specific theorem or later comparison. This dated assertion is retained without independent verification of present-day status.
\subsection{Short English statement}
For \(n-1\) commuting locally nilpotent derivations independent over the polynomial ring, ask whether their joint invariants are generated by a polynomial coordinate. Ring-linear independence is not everywhere pointwise independence, and an arbitrary one-variable invariant ring would be a weaker conclusion.
\subsection{Sources}
\begin{itemize}\item Local UnsolvedMath registry, permanent ID 30000653, code catalogue entry; \emph{Commuting Locally Nilpotent Derivations Conjecture}, labeled commuting derivations conjecture \(\mathrm{CD}(n)\) in the original statement; Oberwolfach Reports Open Problems.\item \emph{Affine Algebraic Geometry (2007)}, Oberwolfach Reports, DOI 10.4171/owr/2007/01, pp. 5--82. \url{https://doi.org/10.4171/owr/2007/01}. Proposed year 2007; no proposer named.\item The same full report is the sole supplied literature source; no separate authoritative coordinate or derivation theorem is supplied.\item Literature triage checked 2026-08-21, status \texttt{open}; record updated 2026-09-16. The \texttt{corrected\_verified} field records label removal and notation normalization. Verify dated claims independently.\end{itemize}
\end{document}
